% ===== FILE: sources/2607.00466-eldr-expert-locality-aware/arxiv/main.tex =====
\documentclass[sigplan, 10pt, nonacm, screen]{acmart}

\renewcommand\footnotetextcopyrightpermission[1]{}
\settopmatter{printfolios=true}
\AtBeginDocument{
  \providecommand\BibTeX{{Bib\TeX}}}

\usepackage{xspace}
\usepackage{booktabs}
\usepackage{tabularx}
\usepackage{xcolor}
\usepackage{kotex}
\usepackage{enumitem}
\usepackage{algorithm}
\usepackage{algpseudocode}
\usepackage{listings}
\usepackage[nameinlink]{cleveref}

\crefname{figure}{fig.}{figs.}
\Crefname{figure}{Fig.}{Figs.}
\crefformat{section}{\S#2#1#3}
\crefmultiformat{section}{\S\S#2#1#3}{ and~#2#1#3}{, #2#1#3}{, and~#2#1#3}

\makeatletter
\newcommand\paragraphnew{\@startsection{paragraph}{4}{\parindent}
  {0pt}
  {-0.4em}
  {\ACM@NRadjust{\@parfont\@adddotafter}}}
\makeatother

\makeatletter
\newcommand{\subsubsectionnew}{
  \def\@toclevel{3}
  \@startsection{subsubsection}{3}{\z@}
  {-.25\baselineskip \@plus -1\p@ \@minus -.1\p@}
  {-3.5\p@}
  {\ACM@NRadjust{\@subsubsecfont\@adddotafter}}
}
\makeatother

\lstset{
  basicstyle=\ttfamily\small,
  columns=fullflexible,
  keepspaces=true,
  frame=single,
  showstringspaces=false
}


\def\sectionautorefname{$\S$}
\def\subsectionautorefname{$\S$}
\def\subsubsectionautorefname{$\S$}


\newcommand{\sys}{\textsc{ELDR}\xspace}

\begin{document}
\pagestyle{plain}
\title{\sys: Expert-Locality-Aware Decode Routing for PD-Disaggregated MoE Serving}

\author{Sangjin Choi}
\authornote{Work done during an internship at Microsoft Research.}
\affiliation{
  \institution{KAIST}
  \city{Daejeon}
  \country{Korea}
}

\author{Sukmin Cho}
\authornotemark[1]
\affiliation{
  \institution{KAIST}
  \city{Daejeon}
  \country{Korea}
}

\author{Yifan Xiong}
\affiliation{
  \institution{Microsoft Research}
  \city{Beijing}
  \country{China}
}

\author{Ziyue Yang}
\authornote{Work partially done while the author was at Microsoft Research.}
\affiliation{
  \institution{Shanghai Xingyunzhili \\ Artificial Intelligence Institute}
  \city{Shanghai}
  \country{China}
}

\author{Youngjin Kwon}
\authornote{Corresponding author.}
\correspondingauthor
\affiliation{
  \institution{KAIST}
  \city{Daejeon}
  \country{Korea}
}

\author{Peng Cheng}
\affiliation{
  \institution{Microsoft Research}
  \city{Redmond}
  \state{WA}
  \country{USA}
}

\input{text/00_abstract}

\maketitle

\input{text/01_introduction}
\input{text/02_background}
\input{text/03_motivation}
\input{text/04_design}
\input{text/05_implementation}
\input{text/06_evaluation}
\input{text/07_related}
\input{text/08_conclusion}

\bibliographystyle{ACM-Reference-Format}
\bibliography{sample-base}

\end{document}



% ===== FILE: sources/2607.00466-eldr-expert-locality-aware/arxiv/text/00_abstract.tex =====
\begin{abstract}
In prefill--decode (PD) disaggregated LLM serving, each request is
assigned to a decode worker after prefill. Existing decode routers
balance only load; for mixture-of-experts (MoE) models this is
incomplete: equally loaded workers can differ in latency, since each
decode step loads the weights of every distinct expert its batch
activates.
We present \sys, an expert-locality-aware decode router for
PD-disaggregated MoE serving. From a request's prefill expert
activations, \sys builds an expert signature predicting the experts it
will activate during generation. Offline, balanced $K$-means partitions
signature space across decode workers; online, locality-band routing
sends each request to the least-loaded worker among those best matching
its signature. A signature cache, co-indexed with the KV cache at
KV-block granularity, keeps signatures exact under prefix caching.
Implemented in vLLM and evaluated on deployments of up to 40 GPUs,
\sys reduces median TPOT by 5.9--13.9\% over the strongest of four
load-balancing baselines across three MoE models and two workloads,
with model outputs unchanged.
\end{abstract}


% ===== FILE: sources/2607.00466-eldr-expert-locality-aware/arxiv/text/01_introduction.tex =====
\section{Introduction}
\label{sec:introduction}

Large Language Model (LLM) serving is moving toward Prefill-Decode (PD)
disaggregation, which runs prompt processing (prefill) and token generation
(decode) on separate worker pools~\cite{distserve}. The phases mismatch: prefill
runs the prompt in parallel and is compute-bound, while decode is sequential and
latency-sensitive, so colocating them lets long prefills stall decode and hurts
both Time-to-First-Token (TTFT) and Time-per-Output-Token (TPOT). An $x$P$y$D
deployment instead provisions $x$ prefill and $y$ decode workers
separately~\cite{dynamo,llmd,deepseekv3,mooncake}, turning routing into the
central serving decision: each request is assigned a prefill worker for its
prompt and, once its KV state is materialized, a decode worker for its
generation. Prefill-side routing has drawn most of the attention, where
cache-affinity policies reuse KV state across prompts. Decode-side routing has
not: existing policies balance load and otherwise treat decode workers as
interchangeable. For dense models they are, since equal-load workers do equal
feed-forward work. For mixture-of-experts (MoE) models they are not.

For MoE models, equal load does not mean equal latency. Decode is
memory-bandwidth bound, and at batched decode its cost is set by the \emph{union}
of distinct experts the batch loads from HBM each step. Here sparsity inverts.
The property that makes MoE cheap, routing each token to only a few
experts~\cite{qwen3,gptoss,gemma4}, fragments the decode batch across
experts and destroys the weight reuse a dense batch enjoys, where one weight load
amortizes over every token. A step pays for every expert any of its tokens
selects, so the union, not the token count, governs latency. The gap is large:
on Qwen3-30B-A3B, growing the active-expert count from 16 to 128 raises MoE-layer
latency $4.7\times$ at fixed batch size, while batch size at fixed active-expert
count barely moves it (\cref{sec:motivation}).
Because the union depends on which requests share a worker, expert
composition is a first-order latency knob that load-based routing
cannot see---a second routing axis: \emph{expert locality}.

Expert locality is exploitable because it is structured.
The MoE gate%
\footnote{We call the MoE router the MoE gate to avoid confusion with the PD router.}
picks each
token's experts from its hidden representation, so requests from the same domain
activate overlapping experts: across tasks, code, math, medical, and legal
prompts exercise distinct expert regions, and multilingual traffic separates by
language. Expert choice is thus correlated across related requests, not just
sparse within a token. Concretely, we observe same-domain decode batches activate
$17$--$21\%$ fewer distinct experts per step than mixed batches on task workloads
(\cref{sec:motivation}), so a router that colocates similar requests shrinks each
worker's per-step union while a load-only router scatters them.

Structure helps only if the router can see it in time. Placement happens at the
prefill-to-decode handoff, before any output token, so decode-time expert choices
are not yet observable. But prefill has already pushed the prompt through the same
gates that will route decode, and the two agree: per-expert prefill and decode
activation correlate at $0.70$ to $0.92$ across three MoE models
(\cref{sec:motivation}). Expert locality is therefore not only structured but
visible at exactly the moment the router must act.

We present \sys, an expert-locality-aware decode router for PD-disaggregated MoE
serving. \sys turns each request's prefill activations into a compact
\emph{expert signature} whose distances predict decode-time expert overlap, so
requests with nearby signatures share experts at decode. It selects the signature
representation by a single criterion, how faithfully signature proximity predicts
that overlap, decoupled from any downstream clustering or routing.

Routing must then satisfy two objectives that need different information.
Locality, which experts colocated requests share, is an aggregate property
visible only across many requests; load, which decode's non-MoE work imposes per
colocated request, is instantaneous and visible only at the decision. Pure
locality overloads popular domains; pure load balancing scatters expert-similar
requests. \sys splits the decision along that information boundary. Offline,
\emph{balanced $K$-means} partitions signature space into one locality region per
decode worker, capturing structure without inheriting workload skew. Online,
\emph{locality-band routing} keeps the workers whose centroids fall within a
similarity band of the request's best match and routes to the least loaded among
them: the band enforces locality, the load choice inside it absorbs live skew,
and one decision serves both.

One case stresses the online path: prefix caching. A cache hit skips
prefill for the shared prefix, so the gate never runs there and the
signature is incomplete---worst exactly for the cache-hit requests
that prefix-aware routing concentrates. \sys therefore keeps an
expert-signature cache co-indexed with the KV cache at block
granularity: each block carries its tokens' expert footprint, and
summing cached and freshly computed blocks recovers the full
signature---coherent across partial hits, evictions, and reuse, with
no recomputation and no change to caching.

We implement \sys in vLLM~\cite{vllm} as a thin layer over an existing
PD-disaggregated stack: prefill-time signature capture, an offline fit,
locality-band routing, and the block-granular signature cache. \sys
changes only which decode worker serves a request, leaving the model,
gate decisions, kernels, and batching untouched, so expert
selections---and hence outputs---are identical to standard top-$k$
gating. Across three MoE models
(Qwen3-30B-A3B~\cite{qwen3}, GPT-OSS-120B~\cite{gptoss},
Gemma-4-26B-A4B~\cite{gemma4}) on task and language workloads,
\sys cuts median TPOT by $7.0$--$13.9\%$ (task) and
$5.9$--$10.0\%$ (language) over the best load-balancing
baseline with a signature cache under $1\%$ of the KV cache.
\sys further generalizes to a 235B deployment under expert parallelism (\cref{sec:eval_generalization}).

\sys makes three contributions:

\begin{itemize}[leftmargin=*, itemsep=0pt, topsep=0pt]
\item \textbf{Expert locality as a predictable routing axis.} Decode latency in
PD-disaggregated MoE serving turns on the distinct experts colocated requests
activate; this locality is structured by domain; and, the enabler, it is readable
from prefill before decode begins.

\item \textbf{Expert-locality-aware decode routing.} \sys builds a
quality-selected expert signature, partitions it across decoders with
balanced $K$-means, and routes with a locality band that balances expert
locality against live load.

\item \textbf{Prefix-cache-coherent signatures.} A block-granular signature cache
co-indexed with the KV cache keeps signatures exact across partial hits, full
hits, and evictions, so \sys composes with prefix caching at negligible cost.
\end{itemize}


% ===== FILE: sources/2607.00466-eldr-expert-locality-aware/arxiv/text/02_background.tex =====
\section{Background}
\label{sec:background}

\subsection{Mixture-of-Experts Large Language Models}
\label{sec:background_moe}

Mixture-of-Experts (MoE) has become a primary architecture for scaling frontier Large Language Models (LLMs) because it decouples model capacity from per-token computation. Instead of applying the same dense feed-forward network (FFN) to every token, an MoE layer replaces the FFN with multiple expert FFNs and an MoE gate that activates a small subset of experts for each token. As a result, the model can contain many more parameters than a dense model while using only a fraction of them during each forward pass. This sparse activation makes MoE
attractive for efficient scaling: it increases model capacity and enables richer
expert specialization without proportionally increasing per-token FLOPs. Recent
architectures push this idea further with fine-grained experts, splitting
FFN capacity across more, smaller experts so that each token can be routed to a
more specialized combination of experts while keeping per-token FLOPs bounded.

\begin{figure*}[!t]
\centering
\includegraphics[width=.88\textwidth]{figure/03_motivation/domain_2x3_layer.pdf}
\vspace{-1em}
\caption{Decode-phase per-expert activation relative to the cross-domain mean,
for three MoE models along task (top) and language (bottom,
WildChat~\cite{wildchat}) domains at each model's most discriminative layer.
Darker is above-average (below-average clipped to white); experts are reordered
per panel into contiguous per-domain blocks. Each domain over-activates a
distinct subset of experts. Task domains:
Code~\cite{humaneval, bigcodebench, ds1000, mbpp},
Math~\cite{gsm8k, math,math500, olympiadbench, aquarat},
Medical~\cite{medqa, pubmedqa, mmlu}, and
Legal~\cite{lexglue}.}
\label{fig:domain}
\end{figure*}


\subsection{Prefill-Decode Disaggregated Serving}
\label{sec:background_pd}

Prefill-Decode (PD) disaggregation~\cite{distserve} is a practical
architecture for latency-sensitive LLM serving. Inference has two
phases: \emph{prefill} processes the prompt in parallel and is
compute-bound, while \emph{decode} generates tokens autoregressively and
is latency-sensitive and memory-bandwidth-bound. Colocating them causes
phase interference---long prefills stall latency-sensitive decode
steps---and forces one resource allocation onto two mismatched
workloads. PD disaggregation separates the phases onto independent
pools, typically an \emph{xPyD} configuration of $x$ prefill and $y$
decode workers, scaling capacity per phase.

This makes routing a central serving decision. A \emph{PD router}
assigns each request a prefill worker for the prompt and a decode worker
for generation, determining where its KV state transfers and which
worker serves its tokens. Existing policies optimize system-level
objectives---cache-aware routers prefer prefillers holding reusable KV
blocks~\cite{mooncake}, load-balancing policies (round-robin,
join-shortest-queue, power-of-two-choices) spread requests---but decide
only \emph{where} a request runs, agnostic to the model-internal expert
activations that set MoE decode's memory-bandwidth cost. This leaves the
activated-expert cost each request imposes on its decode worker
unmodeled.


% ===== FILE: sources/2607.00466-eldr-expert-locality-aware/arxiv/text/03_motivation.tex =====
\section{Motivation}
\label{sec:motivation}

\begin{figure}[!t]
\centering
\includegraphics[width=0.95\columnwidth]{figure/03_motivation/moe_latency.pdf}
\vspace{-1em}
\caption{MoE layer latency scales with active experts, not batch
size (single MoE layer, one MI300X).}
\Description{}
\label{fig:moe_latency}
\end{figure}


\subsection{Active Expert Count Drives MoE Decode Latency}
\label{sec:motivation_active_experts}

MoE decode latency is governed by the number of distinct experts
activated at each decode step. Sparsity reduces computation but
amplifies decode's memory-bandwidth bottleneck by fragmenting the
batch across experts. In a dense model, every token in a decode batch
executes the same FFN, so one weight load amortizes over the whole
batch. In an MoE model, the gate partitions the batch among experts,
and weight reuse occurs only among tokens sharing an expert; as
experts grow more numerous and finer-grained, fragmentation lowers
arithmetic intensity further. MoE decode cost is therefore dominated
by expert weight access: each step fetches the weights of every expert
the batch activates.



To isolate this effect, we benchmark the single-layer MoE expert
computation on one MI300X, sweeping decode batch size and the number
of distinct active experts for Qwen3-30B-A3B, GPT-OSS-120B, and
Gemma-4-26B-A4B. Across all models, latency tracks active experts far
more strongly than batch size (\Cref{fig:moe_latency}): on
Qwen3-30B-A3B, growing active experts from 16 to 128 raises latency
$4.7\times$ at batch size 64, while batch size at a fixed expert count
barely moves it.

This result shows that MoE decode efficiency depends on expert reuse
within a batch. However, exploiting this requires the serving system to know
which requests are likely to activate overlapping experts. We next show
that this overlap is not random. Instead, expert usage is structured by
request domain, making expert locality a predictable property that a serving
system can exploit.

\subsection{Experts Specialize by Domain}

The MoE gate is input-dependent by design. The gate scores experts using
the token's hidden representation and dispatches the token to the top-scoring
experts. Because the hidden representation encodes contextual features, domains
such as task and language are natural axes along which expert usage may differ.

\Cref{fig:domain} shows decode-phase expert activation across three
models: Qwen3-30B-A3B, GPT-OSS-120B, and Gemma-4-26B-A4B. The figure evaluates
two sources of request structure: task and language. The task axis uses a
mixture of code, math, medical, and legal benchmarks. The language axis
uses multilingual WildChat requests, a corpus of one million real-world
user--ChatGPT interactions spanning 68 languages such as English, Chinese,
Russian, and French. In each heatmap, rows correspond to domains and
columns correspond to experts, with color indicating activation
relative to the average for that expert across domains. Experts are sorted by
their dominant domain, so experts most activated by the same domain appear
contiguously.

Across models, expert activation is strongly structured by both task and
language. Code, math, medical, and legal requests over-activate
different expert subsets. The same pattern appears on the multilingual axis:
English, Chinese, Russian, and French requests activate different expert
regions rather than uniformly exercising the full expert pool. This shows that
expert selection is correlated across related requests, not just sparse at the
individual-token level.

This specialization makes active-expert reduction actionable for serving:
same-domain requests are more likely to share experts at decode time, so
co-locating them on a worker yields a smaller per-step active-expert set
than mixed batches. The active-expert count therefore depends on which
requests are placed together, not only on how much work each worker
receives. This structure is a property of the model's gating networks---stable
across requests of the same domain and observable wherever the
gating networks are exercised. To exploit it for placement, a serving
system must identify each request's expert footprint at the moment of
prefill$\to$decode handoff, before any decode token has been generated.

\subsection{Prefill Predicts Decode Expert Activation}
\label{sec:motivation_prefill_decode}

\begin{figure}[t!]
\centering
\includegraphics[width=0.95\columnwidth]{figure/03_motivation/pd_scatter.pdf}
\vspace{-1em}
\caption{Prefill expert activation predicts decode activation. Each
point is one expert (normalized prefill $x$ vs.\ decode $y$, pooled over
domains); points near the diagonal are experts used about equally in
both phases.
}
\label{fig:pd_scatter}
\end{figure}

Expert locality can guide a decode routing policy only if it is visible before
decode begins. In an $x$P$y$D deployment, a request first runs on a prefill worker
and is then assigned to one of the decode workers, where it will be batched with
other active requests. The decode routing policy must make this placement
decision before any output tokens are generated, so the request's decode-phase
expert usage is not yet observable. The prefill phase, however, has already
processed the prompt through the same MoE layers, exposing the expert choices
made by the model's gating networks.

\Cref{fig:pd_scatter} compares prefill and decode expert activation across
the task and language domains studied above. Each point corresponds to one
expert: the $x$-axis is its normalized activation during prefill, and the $y$-axis is
its normalized activation during decode. Points near the diagonal indicate that
experts heavily used while reading the prompt are also heavily used while
generating the response. The correlation is strong for Qwen3-30B-A3B and
Gemma-4-26B-A4B, and substantial for GPT-OSS-120B.

This correlation makes decode expert locality available before decode begins.
Although the exact decode tokens are unknown at placement time, the prefill
phase has already exposed a request-specific expert footprint that predicts the
experts likely to be reused during generation. Thus, in an $x$P$y$D deployment, the
system observes an expert-locality signal precisely at the boundary where the
request leaves the prefill worker and must be assigned to a decode worker.

\subsection{Opportunity: Expert-Locality-Aware Decode Routing}
\label{sec:motivation_opportunity}


In PD-disaggregated MoE serving, decode-step latency is governed not
only by per-worker load but by request composition: requests that
dispatch to overlapping experts activate fewer distinct experts per
step than an equally sized batch of unrelated requests.
\Cref{fig:opp_sweep} quantifies this: same-domain batches
activate 17--21\% (task) and 3--10\% (language) fewer experts per
step than mixed-domain batches across batch sizes 32--128.

Expert locality offers a second routing axis for
PD\--dis\-ag\-gre\-gat\-ed MoE serving.
Existing decode routing policies
balance per-worker load but ignore expert overlap. Placing
expert-similar requests on the same decoder shrinks the per-step
active-expert set; placing dissimilar requests together expands
it. The signal that enables overlap-aware routing is observable
at the prefill$\,\to\,$decode handoff: prefill expert activations
predict decode-time expert usage. We call the per-request
representation of this signal the \emph{expert signature}.


\begin{figure}[t]
\centering
\includegraphics[width=0.95\columnwidth]{figure/03_motivation/active_experts_sweep.pdf}
\vspace{-1em}
\caption{Same-domain batches (blue) activate fewer experts per decode
step than mixed-domain (orange),
for task (top) and language (bottom) across three MoE models.}
\label{fig:opp_sweep}
\end{figure}



\subsection{Challenges}
\label{sec:motivation_challenges}

Turning this opportunity into a working routing policy raises three
challenges.

\subsubsectionnew{Designing the expert signature}
The expert signature predicts a request's decode-time expert
usage from its prefill activations, so that requests with
similar predicted usage can be colocated to share experts.
It is the building block of \sys's design: every routing
decision derives from it, and routing quality is bounded above
by signature quality.
The challenge is to turn raw prefill data into a vector space
suited to clustering---one in which proximity between two
signatures reflects how much their requests overlap in
decode-time expert activation. Only such a space lets a
clustering group requests by genuine expert affinity rather
than by spurious similarity. The design space is large: raw
activation counts, gate logits, layer masks, and reweighting
schemes all yield candidate representations, each inducing a
different geometry. Selecting among them requires a criterion
that measures how faithfully signature proximity predicts
decode-time expert overlap, independent of any particular
clustering or routing policy built on top.

\subsubsectionnew{Reconciling locality and load}
\label{sec:challenge_load}
Locality alone is not enough. Non-MoE compute---attention, QKV
projection, normalization---scales with the number of active
requests on a worker, not with how many distinct experts they
share. A locality-only policy concentrates traffic on the
workers whose profiles match the dominant request domains; in
WildChat, the top two languages alone account for $\sim$75\%
of requests (\Cref{fig:wildchat}). Overloaded workers
see inflated non-MoE work and rising tail TPOT. The opposite
extreme---a balance-only policy like round-robin or pure
JSQ---keeps queues even but scatters expert-similar requests,
erasing the locality benefit. Both objectives must hold at once.


\begin{figure}[t]
\centering
\includegraphics[width=0.95\columnwidth]{figure/03_motivation/wc_langdist.pdf}
\vspace{-1em}
\caption{WildChat~\cite{wildchat} request volume is heavily skewed: English and
Chinese alone are $\sim$75\% of requests.
}
\label{fig:wildchat}
\end{figure}

Combining the two is hard because they draw on different information.
Locality is purely \emph{aggregate}: which signatures group together is
visible only across many requests, never from one. Balancing load needs
that aggregate view too---structural skew, where some clusters carry far
more traffic than others, lets the router anticipate the distribution
instead of reacting request by request---but also an \emph{instantaneous}
view: runtime variance from Poisson arrivals, variable output lengths,
and prior routing decisions that no aggregate can predict, so the router
can track each worker's live load at the moment of decision. No single
mechanism provides both: aggregate-only is blind to live load,
instantaneous-only discards the cross-request structure. The challenge is
to bring both views into the routing pipeline.

\subsubsectionnew{Maintaining signature coherence across prefix
cache hits}
\label{sec:challenge_prefix}

The expert signature is produced during prefill, but prefill is
often partially or fully skipped. With a prefix cache hit, the
system reuses cached KV states for the shared prompt prefix
without re-running the gating networks on the cached tokens.
The expert footprint for those tokens is therefore absent from
this request's prefill---even though it was produced (and
discarded) when an earlier request first populated the cache.
Without a mechanism to recover this footprint, the signature
for a cache-hit request is incomplete and misroutes the request.

A natural answer is to cache the per-request signature.
This works on full prompt matches but breaks on \emph{partial hits}:
a request can share its leading KV blocks with a previous
one and diverge on the rest, and no whole-prompt signature
describes this kind of overlap. KV blocks are also dynamically
evicted, and recycled across requests, so any signature
mechanism must remain coherent under these transitions. The challenge
is to keep the signature coherent with the cache across these
states---partial hits, evictions, and reuse---without disabling
the cache or re-running prefills.


% ===== FILE: sources/2607.00466-eldr-expert-locality-aware/arxiv/text/04_design.tex =====
\section{Design}
\label{sec:design}

\begin{figure}[t]
\centering
\includegraphics[width=\columnwidth]{figure/04_design/fig_overview.pdf}
\vspace{-2.0em}
\caption{\sys architecture: offline fitting of one centroid
per decode worker over expert signatures, then online routing
at the prefill--decode handoff by signature similarity,
subject to load.}
\label{fig:eldr_overview}
\end{figure}

We present \sys, an \textbf{E}xpert-\textbf{L}ocality-aware
\textbf{D}ecode \textbf{R}outing policy for PD-disaggregated MoE
serving. Instead of placing decode requests by load alone, \sys
uses each request's prefill expert activation to choose a decode
worker where it is likely to share experts with other
requests. This makes expert locality a practical routing signal
for reducing TPOT.

\subsection{Overview}
\label{sec:design_overview}

\Cref{fig:eldr_overview} shows the \sys architecture.
\sys integrates into an existing PD-disaggregated serving stack
as a thin layer: a request router in front of the prefill--decode
workers and a prefill-time hook that records per-block \emph{(1) expert
signatures} alongside the KV cache, with offline routing state
loaded once at startup. The model, kernels, batching, and the
prefill/decode engines are left unmodified. We implement and
evaluate \sys on vLLM~\cite{vllm}, and the design ports directly
to other stacks such as SGLang~\cite{sglang}. \sys runs in two
stages: an \emph{(2) offline stage} that groups requests by their
prefill-time signatures to reduce expert activation, and an
\emph{(3) online stage} that uses this profile to route each request
at the prefill--decode handoff.

\paragraphnew{Expert Signature}
\sys builds on the observation that a request's prefill expert
activations are highly correlated with its decode expert
activations (\cref{sec:motivation_prefill_decode}). To
exploit this, we design a per-request representation, the
\emph{expert signature} (\cref{sec:design_signature}), that
summarizes a request's prefill-time expert activations and predicts
which experts it will use at decode time. The signature lets us
estimate how much two requests' decode-time expert usage overlaps
directly from prefill information, via the cosine similarity between
their signatures. \sys therefore groups requests by signature
similarity, so that the distinct experts activated across a group
are fewer than under naive, similarity-agnostic routing.

\paragraphnew{Offline}
Before serving, \sys runs a small calibration set through prefill and
collects each request's expert signature. \sys then clusters these
signatures (\cref{sec:design_clustering}) along two axes:
locality, so signatures grouped on one worker share experts, and
balance, so clusters are evenly sized across decode workers (one centroid
per worker). Re-fitting is cheap: the balanced $K$-means clustering over the
captured signatures completes in under $10$\,s on CPU and updates only
the router's centroid table, leaving the model and decode workers
running. A PD scale-up or scale-down, or a shifted workload, then needs
only this re-fit at the new decode-worker count, reassigning one
centroid per worker. We use a single offline fit in this work.

\paragraphnew{Online}
During serving, \sys leaves the standard PD pipeline intact and
adds two lightweight steps: signature capture during prefill and
a routing decision at the prefill--decode handoff:

\begin{enumerate}[leftmargin=*, topsep=0pt]
\item The router sends the request to a prefill worker using
prefix-aware routing.

\item The prefill worker runs prefill while \sys captures the
request's \emph{expert signature}
(\cref{sec:design_signature}) in a signature cache
co-indexed with the KV cache
(\cref{sec:design_prefix_cache}).

\item At handoff, the request's expert signature---summed over
its expert signature cache blocks---is returned to the router
with the KV-transfer metadata.

\item The router invokes \sys's \emph{locality-band routing}
(\cref{sec:design_route}): it compares the signature
against the decoder centroids and routes to the least-loaded
worker within the locality band of the closest centroid.

\item The selected decoder pulls the KV cache and generates on the
unchanged decode engine; co-locating similar-signature requests shrinks
the per-step active-expert union, reducing decode latency.
\end{enumerate}

\subsection{Expert Signature}
\label{sec:design_signature}

\sys summarizes each request's prefill-time expert activations into
a compact per-request representation, the \textbf{expert signature}.
The raw material is cheap to collect at prefill: a per-layer
histogram of how many tokens each expert receives, available as a
by-product of routing. The signature is designed so that the
distance between two signatures predicts how much the two requests
overlap in their decode-time expert usage---requests with nearby
signatures activate similar experts at decode, requests with distant
signatures activate different ones. The clustering layer
(\cref{sec:design_clustering}) then groups requests whose signatures
are close, so each decode worker loads few distinct experts per step.

Turning the raw histogram into a good signature still takes care.
The per-layer counts are the right raw signal, but two effects keep
raw count distance from tracking decode-time overlap as closely as
it could: a few generalist experts fire on nearly every request,
inflating the norm and masking the rare specialists that distinguish
workloads, and many layers carry little signal that separates
requests. The signature therefore weights the counts and selects
which layers to keep, on top of what to count. We proceed in three
steps: we make the design goal precise (\cref{sec:sig_metric})---a
signature is good when its distances rank request pairs the same way
their decode-time expert overlap does---then construct the signature
from three principles justified against that goal
(\cref{sec:sig_construction}), and validate each choice empirically
(\cref{sec:sig_validation}).

\subsubsectionnew{Goal of Signature Design}
\label{sec:sig_metric}

The signature design aims to make decode-time expert overlap
predictable from prefill alone, so clustering on signatures groups
requests that load few distinct experts per step and thereby bounds
decode latency. A signature meets this goal when requests with nearby
signatures activate similar experts at decode time, and distant
signatures activate different ones.

To quantify it, we encode each calibration request $i$'s decode-time
behavior as a per-step \emph{activation probability}
$p_i(\ell, e) \in [0, 1]$: the fraction of $i$'s decode steps in which
the gate at layer $\ell$ selected expert $e$. Flattened across all
$(\ell, e)$ pairs and L2-normalized, $p_i \in \mathbb{R}^{LE}$ is the
decode-time pattern the signature must predict. Distance in $p$-space
is the natural quantification of decode-time expert overlap. Signature
quality is then the rank correlation between signature pair-distance and
decode-pattern pair-distance, over random calibration pairs:
\begin{equation}
\rho \;=\; \mathrm{Spearman}\big(\,
  \mathrm{cos\text{-}dist}(s_i, s_j),\;\mathrm{cos\text{-}dist}(p_i, p_j)
\,\big).
\label{eq:rho}
\end{equation}

We use Spearman (rank) correlation rather than Pearson: signature
distance lives in $\mathbb{R}^d$ while decode-pattern distance lives in
$\mathbb{R}^{LE}$, so their numerical values are not commensurable.
Only the ordering of pairs is, and ordering is exactly what the
clustering layer consumes when deciding which requests to colocate.
Rank agreement also captures the goal property: high $\rho$ means nearby
signatures correspond to similar decode-time expert activation, and
distant signatures to different activation.

\subsubsectionnew{Signature Design}
\label{sec:sig_construction}
Three principles guide $s_r$:
\begin{itemize}[leftmargin=*, itemsep=2pt]
\item \textbf{Discrete, not continuous.} The target $p_i(\ell,e)$ is the
normalized count of top-$k$ selections at each $(\ell,e)$ over decode, so
the signature should measure the same quantity. Only the discrete
top-$k$ loads experts---sub-threshold scores are computed but never fetch
one---so a discrete prefill count is a sparse, same-form estimator of
which experts fire. Continuous gate scores instead spread softmax mass
over all $E$ experts, assigning magnitude to experts the top-$k$ never
loads and smoothing away the contrast that distinguishes requests.

\item \textbf{Downweight common experts.} A handful of generalist
experts fire on almost every prefill. They carry little information
about request identity, yet their large counts drown out the rare
specialist experts that actually distinguish workloads. We multiply
each $(\ell, e)$ count by its inverse document frequency (IDF) over
the calibration corpus, which shrinks the common experts and
amplifies the rare ones. This reweighting works at the granularity of
$(\ell, e)$ cells, denoising the signature \emph{before} we choose a
layer mask.

\item \textbf{Keep informative layers.} Layers contribute unequally:
some carry routing signal that maps prefill to decode, others add
dimensions that don't separate requests and dilute the signature
direction. The mask $\mathcal{S}\subseteq\{1,\ldots,L\}$ is the subset
we keep, chosen per deployment to maximize $\rho$
(Eq.~\ref{eq:rho}) on calibration data.
\end{itemize}

Four steps build $s_r$. The IDF table and layer mask are fit once
offline on calibration data; counting, IDF reweighting, layer masking,
and normalization run online at the prefill$\,\to\,$decode handoff.

\begin{enumerate}[leftmargin=*]
\item \emph{Count.} At each layer $\ell$, count how many prefill tokens
are routed to each expert, giving $c_r(\ell)\in\mathbb{N}^{E}$.

\item \emph{Downweight common experts.} Multiply each cell by its
inverse document frequency,
$w(\ell,e)=\log\!\big((|\mathcal{C}|+1)/(\mathrm{df}(\ell,e)+1)\big)$,
where $\mathrm{df}(\ell,e)$ counts calibration requests
$\mathcal{C}$ in which expert~$e$ fires at least once at layer~$\ell$.
The reweighted count is
$\tilde c_r(\ell,e)=c_r(\ell,e)\cdot w(\ell,e)$.

\item \emph{Select layers.} The mask~$\mathcal{S}$ is fit offline by
greedy layer selection on $\rho$: starting from an empty mask,
the layer whose inclusion most increases cumulative~$\rho$ is appended
at each step, producing an ordering of all $L$ layers. We keep the
first $N^{*}$ layers in this ordering, where $N^{*}$ is the size at
which cumulative~$\rho$ peaks (the~$\bigstar$ in~\Cref{fig:sig_layers}).
At runtime, the per-request reweighted counts on $\mathcal{S}$ stack to
$x_r=\big[\,\tilde c_r(\ell)\,\big]_{\ell\in\mathcal{S}}\in\mathbb{R}^{N^{*}\cdot E}$.

\item \emph{Normalize.} Divide by the vector's length so similarity
reflects the shape of expert usage, not the request's token count:
\begin{equation}
s_r \;=\; x_r \,/\, \lVert x_r\rVert_2.
\label{eq:sig}
\end{equation}
\end{enumerate}

The signature $s_r$ is the building block consumed by the decode
clustering and routing layers (\cref{sec:design_routing}).

\begin{figure}[t]
\centering
\includegraphics[width=0.95\columnwidth]{figure/04_design/sig_transform.pdf}
\vspace{-1em}
\caption{Signature quality~$\rho$ (Eq.~\ref{eq:rho}) for six candidate
transformations~$T$. Bars are the mean across six cells
(3 models~$\times$ 2 workloads); whiskers span the per-cell min/max.}
\label{fig:sig_transform}
\end{figure}


\subsubsectionnew{Validation}
\label{sec:sig_validation}
We validate each design choice by its effect on $\rho$
(Eq.~\ref{eq:rho}), measured on 1{,}000 calibration requests per cell
across three models (Qwen3-30B-A3B, GPT-OSS-120B, Gemma-4-26B-A4B) and
two workloads (task, language).

\paragraphnew{Discrete beats continuous}
\Cref{fig:sig_transform} compares six choices of $T$ on the full
layer mask: count and count$\cdot$idf (Steps~1--2 of
\S\ref{sec:sig_construction}), plus four ablations---$\sqrt{\text{count}}$;
\emph{gate prob} and \emph{gate logit}, the summed softmax and raw
gate scores over prefill tokens; and \emph{binary}, the indicator that
an expert appears in the top-$k$ at least once.
The three count-based variants beat both continuous ones by
\textgreater 0.035 in mean~$\rho$: continuous scores put mass on experts
the decode top-$k$ never loads, and that mass cannot appear in~$p_i$.

\paragraphnew{{IDF denoises}}
Among the discrete variants, count$\cdot$idf raises mean~$\rho$ over
plain count by 1.5pt and its worst case ($\min_c \rho$) by 4.9pt
($0.63\rightarrow0.68$ on GPT-OSS language), with no loss on the best
case. The gain concentrates on GPT-OSS, where a few generalist experts
dominate the raw count and inflate the $\ell_2$ norm; IDF down-weights
them and surfaces the rare experts that distinguish workloads (+8.9pt on
GPT-OSS task).

\paragraphnew{Layer mask compresses}
Running greedy layer selection on count$\cdot$idf signatures,
\Cref{fig:sig_layers} plots cumulative~$\rho$ against layers kept.
Every cell peaks at $N^{*}<L$, exceeding the all-layer $\rho(L)$ by
0.005 to 0.032: extra layers add dimensions that don't separate requests
and dilute the signature. \sys's offline fit uses this per-cell
peak~$N^{*}$.

\begin{figure}[t]
\centering
\includegraphics[width=0.95\columnwidth]{figure/04_design/sig_layers.pdf}
\vspace{-1em}
\caption{Cumulative~$\rho$ (Eq.~\ref{eq:rho}) versus the number of
layers kept under greedy layer selection. One panel per model;
task (blue) and language (orange) shown separately. The star
marks the peak~$N^{*}$ chosen by \sys{}'s offline fit.}
\label{fig:sig_layers}
\end{figure}


\subsection{Decode Clustering and Routing}
\label{sec:design_routing}

Given an expert signature, \sys routes the request to one of
$K$ decode workers under two objectives. \emph{1) Locality}---co-located
requests should share experts---is what the signature is built to
enable. \emph{2) Load balance} is the constraint routing cannot ignore:
non-MoE compute (attention, QKV projection, normalization) scales with
a worker's batch size, so an overloaded worker spikes tail TPOT
(\cref{sec:challenge_load}).
\sys splits these objectives across two stages by the information each
needs. Locality is visible offline from aggregate calibration data, so
the offline stage (\cref{sec:design_clustering}) clusters signatures
into $K$ locality regions of equal calibration volume, one per worker.
Load depends on the instantaneous per-worker state at the
prefill$\,\to\,$decode handoff, so the online stage
(\cref{sec:design_route}) routes each request within its region toward
the less-loaded worker.


\subsubsectionnew{Offline clustering}
\label{sec:design_clustering}

\sys produces one centroid per decode worker by clustering the
calibration signatures into $K$ groups such that signatures within a
group are close (locality) and groups are balanced in size.
$K$-means~\cite{kmeans} is the standard algorithm for locality-driven
clustering: it minimizes the mean cosine distance between each signature
and its assigned centroid. This objective optimizes locality but has no
size penalty---a single centroid in a tight, dense region serves all
those points cheaply because they are close to it, while a centroid
covering a sparse region serves few points but at similar per-point
cost. On skewed data, this produces uneven cluster sizes, and since each
cluster maps to one decode worker, the imbalance translates directly to
uneven decoder load at runtime, causing tail latency spikes. The missing
piece is a size constraint: \sys uses \textbf{Hungarian-balanced
$K$-means}~\cite{balancedkmeans}, which replaces nearest-neighbor
assignment with a globally optimal assignment: each centroid takes at most
$\lceil N/K \rceil$ points, and the assignment minimizing total cosine
distance is found by the Hungarian algorithm. Online routing
(\cref{sec:design_route}) uses cosine similarity to the resulting $K$
unit centroids to identify candidate decoders.
Because $K$ is the number of decode workers, scaling the decode pool up
or down can be absorbed by a re-cluster at the new $K$---a cheap offline
re-fit (under $10$\,s, \S\ref{sec:overhead}) that needs no
re-deployment; the same re-fit could track workload drift.

\Cref{fig:cluster_pca} projects the calibration signatures onto
their first two principal components and overlays the balanced $K$-means
centroids. Task domains (Code/Math/Medical/Legal) and WildChat languages
(English/Chinese/Russian/French) occupy distinct regions across all
three models---the signature space has genuine semantic structure for
the clustering to exploit. The centroids spread across the distinct
regions rather than crowding onto the densest one, which is exactly the
property the balance constraint buys: equal cluster size is only
achievable when centroids cover all the distinct regions, not when they
collapse onto the most populated one. Some centroids consequently land
in sparser regions of the projection---vanilla $K$-means would draw them
onto the densest modes, but the balance constraint pulls them outward to
absorb their $\approx\!N/K$ share.

\begin{figure}[t]
\centering
\includegraphics[width=0.95\columnwidth]{figure/04_design/cluster_pca.pdf}
\vspace{-1em}
\caption{PCA of calibration signatures with Hungarian-balanced centroids
($K{=}8$ for legibility), colored by task domain (top) or top-4
WildChat~\cite{wildchat} language (bottom).}
\label{fig:cluster_pca}
\end{figure}

\subsubsectionnew{Online routing}
\label{sec:design_route}

Pure top-1 routing---sending each request to its single nearest
centroid---spikes tail latency because it ignores instantaneous load:
a burst of requests near one centroid creates load imbalance at that
worker while others sit idle. \sys's online routing rule,
\textbf{locality-band routing}, favors locality when the signature
singles out one worker, and falls back to load-aware balancing when
several workers are nearly tied.

For each request, \sys computes cosine similarity between the
request's expert signature and each of the $K$ centroids; let
$s^* = \max_k s_k$ denote the highest similarity.
Among workers with $s_k \ge s^* - \tau$ (the \emph{locality band}), \sys
selects the one with the smallest load---the number of in-flight decode
requests on that worker, tracked at the proxy and updated as requests are
dispatched and complete. The parameter $\tau \in [0, 1]$
interpolates between pure top-1 routing ($\tau{=}0$) and pure
shortest-queue routing ($\tau{=}1$).

The locality band adapts to signature confidence: a confident
signature that strongly favors one centroid puts few workers in the
band (locality preserved), while an ambiguous signature near several
centroids puts more workers in the band (load balance applied). \sys
uses $\tau{=}0.1$ throughout the paper;~\cref{sec:design-validation}
examines this choice.


\subsection{Prefix Cache Coherence}
\label{sec:design_prefix_cache}



Prefix caching reuses KV blocks across requests that share a prompt prefix,
skipping the prefill of the shared portion entirely. This creates a problem
for \sys. Routing requires the full prompt's expert signature, but the
expert activations for a cached prefix were produced by an earlier request
and are not computed at this request's prefill.

A naive approach is to cache each request's expert signature and reuse it
when an identical prompt arrives. But a request can interact with the
prefix cache in several ways: a prompt may miss the cache entirely, hit
on some leading blocks and compute the rest (a partial hit), or hit on
the entire prompt and skip prefill (a full hit); a cached block may also
be evicted and reused for a different request. A per-request signature
cache only matches on full hits, and any mismatch degrades the expert
locality routing is meant to exploit.


\begin{figure}[t]
\centering
\includegraphics[width=0.95\columnwidth]{figure/04_design/fig_prefix_coherence.pdf}
\vspace{-1em}
\caption{\sys stores expert signatures at KV cache block granularity: the
signature cache is co-indexed with KV cache.}
\label{fig:prefix_cache_coherence}
\end{figure}

The prefix cache manages KV blocks one at a time, so \sys tracks the expert
signature at the same granularity: one row per KV block, indexed by the same
block id. As prefill writes a token's K/V into a block, the top-$k$ experts
selected for that token at each MoE layer are written to the matching signature
row. A request's KV cache spans a block set $\mathcal{B}(r)$, so its signature
is the sum over those blocks, $s_r = \sum_{b \in \mathcal{B}(r)} \mathrm{sig}[b]$,
exact regardless of which request populated each block: on a partial hit
(\Cref{fig:prefix_cache_coherence}), cached blocks contribute rows from
the earlier request that first prefilled them and fresh blocks contribute this
request's rows, and the sum matches a cold prefill of the same prompt. The
signature cache inherits the KV cache's lifecycle---a reused block resets its
signature, a cached prefix carries its signature, and an eviction reclaims both
stores together---so \sys keeps no prefix-tree, eviction, or sharing state of
its own.

\sys's signature cache is a preallocated GPU tensor with one slot per
KV-cache block, using one int8 per (block, MoE layer, expert); its total
size is bounded by $\text{num\_gpu\_blocks} \cdot L \cdot E$ bytes. For
the models we evaluate ($L \le 48$, $E{=}128$), each block's row is at
most $6$\,KiB, keeping the whole cache well under $1\%$ of the KV cache
it shadows. Each forward accumulates the new tokens' expert picks into
the matching rows in a single GPU pass; retrieving a request's signature
sums its blocks' rows into its $[L,E]$ signature, batched into the
post-step device-to-host copy. Indexed by block id and tied to the KV
cache's block lifecycle, the cache adds no allocator, no eviction state,
and no synchronization beyond what the KV cache already provides.


% ===== FILE: sources/2607.00466-eldr-expert-locality-aware/arxiv/text/05_implementation.tex =====
\section{Implementation}
\label{sec:implementation}

\sys is implemented on top of vLLM~\cite{vllm} as three components: a routing
proxy that dispatches each request to a decoder by signature locality,
engine-side hooks that emit per-request expert signatures during prefill, and a
one-shot offline pipeline that clusters captured signatures into per-decoder
centroids---roughly $2{,}000$ lines of new Python.

\paragraphnew{Signature path}
A prefill hook accumulates per-layer per-expert activation counts at KV-block
granularity, stored in an expert signature cache keyed by block id. At prefill
completion, the prompt's per-block signatures are summed and piggybacked onto
the existing prefill--decode handoff---no new RPC channel, and no per-request
state in the proxy.

\paragraphnew{Per-decoder EP load balancing}
Following METRO~\cite{metro}---which balances expert-parallel (EP) ranks by
\emph{activated-expert} load rather than token count---we use the calibration
trace to estimate each cluster's expected per-expert decode activation, then for
every (decoder, layer) greedily assign the next most-active expert to the
lightest-loaded rank. Since each decoder serves a distinct cluster, a single
global EP layout would balance one cluster at the others' expense; per-decoder
placement matches each decoder to the workload it sees while holding the
per-rank expert count fixed (uniform weight memory).

\paragraphnew{Offline fit}
A single script turns the calibration trace into the artifacts loaded at
startup: a routing artifact (layer mask, IDF weights, and $K$
Hungarian-balanced centroids in the masked, IDF-weighted, L2-normalized space)
and, for multi-GPU expert parallelism, an expert placement file. It runs once
per (model, dataset, $K$); recalibration is needed only when the model, dataset,
or prefill--decode topology changes.


% ===== FILE: sources/2607.00466-eldr-expert-locality-aware/arxiv/text/06_evaluation.tex =====
\section{Evaluation}
\label{sec:evaluation}

\begin{figure*}[t!]
\centering
\begin{minipage}[t]{0.49\textwidth}
\centering
\includegraphics[width=\linewidth]{figure/06_evaluation/fig_main_task.pdf}
\vspace{-1.5em}
\caption{TPOT (median, p99) and median TTFT vs request rate on the
\textbf{task} workload at 8P16D.}
\label{fig:main_task}
\end{minipage}\hfill
\begin{minipage}[t]{0.49\textwidth}
\centering
\includegraphics[width=\linewidth]{figure/06_evaluation/fig_main_lang.pdf}
\vspace{-1.5em}
\caption{TPOT (median, p99) and median TTFT vs request rate on the
\textbf{language} workload at 8P16D.}
\label{fig:main_language}
\end{minipage}
\end{figure*}

\paragraphnew{Testbed} We evaluate \sys on a 5-node cluster of
AMD MI300X GPUs (8 GPUs per node, 192 GB HBM each), interconnected
by 400 Gbps NDR InfiniBand with 8 ConnectX-7 HCAs per node (one
per GPU). All experiments run on vLLM~\cite{vllm} 0.21.0rc1 /
ROCm 7.2, with prefill--decode disaggregation served by vLLM's
NIXL~\cite{nixl} connector.

\paragraphnew{Models} We exercise four open MoE models. Three are served at
$TP{=}1$: \textbf{Qwen3-30B-A3B}~\cite{qwen3} (128 experts, top-8, bf16),
\textbf{GPT-OSS-120B}~\cite{gptoss} (128 experts, top-4, mxfp4), and
\textbf{Gemma-4-26B-A4B}~\cite{gemma4} (128 experts, top-8, bf16).
For the large-MoE scaling study, \textbf{Qwen3-235B-A22B}~\cite{qwen3}
(128 experts, top-8, bf16) is served with $TP{=}4$ / $EP{=}4$ per instance.

\paragraphnew{Datasets} Each model is evaluated on two workloads.
\textbf{Task} is an 11,668-prompt mix spanning four domains drawn
from public benchmarks, with naturally unequal domain shares
reflecting source-dataset sizes: 3,600 \emph{legal}~\cite{lexglue},
2,779 \emph{code}~\cite{humaneval, bigcodebench, ds1000, mbpp},
2,744 \emph{math}~\cite{gsm8k,math,math500,olympiadbench, aquarat},
and 2,545 \emph{medical}~\cite{medqa, pubmedqa, mmlu}\footnote{We use the
professional-medicine subset of MMLU~\cite{mmlu}.} --- a 1.41$\times$
largest/smallest ratio that stresses load balance when
routing by signature. \textbf{Language} is a 14,000-prompt subset
of WildChat~\cite{wildchat} that inherits the heavy language skew
shown in~\Cref{fig:wildchat}: English, Chinese, Russian, and
French account for 87.6\% of the volume.
For each dataset, the 1,000-prompt
offline calibration split is disjoint from the serving evaluation
set.

\paragraphnew{Topologies and workload}
Unless noted, the main evaluation uses an 8-prefiller / 16-decoder
(8P16D) topology, with each instance at $TP{=}1$ so the 24 GPUs
straddle 3 nodes. We additionally sweep 8P8D and 8P24D for topology
generalization, and 2P8D~$\times$~($TP{=}4$, $EP{=}4$) for the
235B run, the largest deployment in this study at 40 GPUs spanning
all five nodes of the cluster. Requests are Poisson-distributed at
offered rates from 20 to 100 qps (24--56 qps for 235B), each rate
held for 120s with a 30s warmup. We cap output at 512 tokens and
enforce \texttt{ignore\_eos}.

\paragraphnew{Routers and baselines} Our serving stack follows the prefill/decode
policy split adopted by vLLM~\cite{vllm}, SGLang~\cite{sglang}, and NVIDIA
Dynamo~\cite{dynamo}: cache- and hash-affinity policies on the prefill side,
where the prefiller holds cross-request KV state, and load-balance policies on
the decode side, where each request receives a fresh per-request KV transfer
and the decoder holds no cross-request state. We adopt PrefixHash---from the
affinity family---at our prefill, with a least-loaded fallback when the
prefix-matched prefiller is above a $1.25\times$ average-load threshold. On the
decode side we compare \sys against four load-balance baselines---Random,
Round-Robin (RR), Join-Shortest-Queue (JSQ)~\cite{jsq}, and
Power-of-Two-Choices (P2C)~\cite{p2c}---and Domain, a na\"ive
locality-aware baseline that mirrors \sys but routes on a given oracle
domain label instead of the expert signature: it splits the decoders
among the domains in proportion to the calibration domain mix, and sends
each request to the decoders assigned to its domain, load-balanced among
them by JSQ. Within each cell, all six routers share the same prefill
policy and decoder pool, so a cell isolates the decoder routing decision.

\paragraphnew{\sys configuration} All experiments use the design from~\cref{sec:design}: count$\cdot$idf signatures restricted to a
greedy $\rho$-selected layer subset, Hungarian-balanced $K$-means
with $K{=}D$ (one cluster per decoder), and $\tau{=}0.1$
locality-band routing. The 235B run additionally enables per-cluster
per-layer expert-rank permutation for EP load balancing. Signature
capture takes 4--15 minutes per (model, dataset) for the 1,000
calibration prompts; the offline fit (greedy mask selection plus
balanced $K$-means) completes in under 10 seconds.

\subsection{Main Results}
\sys reduces median TPOT across all three models and both workloads,
with the size of the win tracking how separable the workload's expert
activations are. The two baseline families isolate \sys's two
ingredients: the load balancers add load-awareness without locality,
Domain adds static locality without live load, and \sys combines finer
locality with live load. We report the mean across the five request
rates; the improvement holds across the whole sweep rather than at a
single operating point.

On the task workload (\Cref{fig:main_task}), where prompts fall
into well-defined domains (code, math, medical, legal), \sys reduces
median TPOT by $7.0$--$13.9\%$ and tail TPOT by $3.4$--$6.0\%$ over the
best load balancer, and sits below every one at every rate. Domain is a
stronger baseline here---its labels align with the model's expert
clusters, so a static per-domain partition concentrates each block's
expert working set and itself reduces median TPOT by $6.8$--$9.7\%$ over
the load balancers---yet \sys still beats it, by $1.4$--$6.9\%$ on
median and $1.6$--$4.5\%$ on tail TPOT, because its $K{=}16$ signature
clusters resolve locality more finely than four labels and its
$\tau$-band spills load across cluster boundaries that Domain's hard
partition forbids. \sys's median TTFT tracks the baselines and is lower
near saturation, where its faster decoders relieve prefill
back-pressure.

On the language workload (\Cref{fig:main_language}), where domain
boundaries are softer, \sys reduces median TPOT by $5.9$--$10.0\%$ over
the best load balancer; mean tail TPOT reduces by $6.2\%$ on
Qwen3-30B-A3B and regresses by $1.5\%$ on GPT-OSS-120B and $0.2\%$ on
Gemma-4-26B-A4B, though at the per-cell peak all three reduce ($9.6\%$,
$1.0\%$, $5.6\%$). Domain collapses here: a language label is a coarse
proxy for expert activation---\Cref{fig:cluster_pca} shows several
of \sys's centroids falling within a single language, so one static
block per language mixes unrelated signature clusters---and its skewed
mix over-subscribes one block, so Domain only matches the load balancers
on median TPOT (within $3\%$ of RR) and regresses tail TPOT by up to
$6.1\%$ as the hot block saturates. \sys beats Domain by $5.7$--$9.1\%$
on median and $7.0$--$9.5\%$ on tail TPOT: its finer signature
clusters capture the intra-language sub-structure that the
coarse labels miss (\Cref{fig:cluster_pca}),
and its $\tau$-band routing rebalances live load that the static
partition cannot. Median TTFT again tracks the baselines.

\subsection{Overhead Analysis}
\label{sec:overhead}
\sys's overhead has a one-time offline part and a per-request serving
part. Offline, the router is built once per deployment: a
signature-capture pass over the $1{,}000$ calibration prompts, then the
fit---greedy mask plus balanced $K$-means---which completes in under
$10$~s on CPU and is cheap to re-run on configuration changes. At
serving time (\Cref{tab:overhead}), \sys adds $0.86$~ms per
request, $1.2\%$ of the $69$~ms median TTFT, dominated by per-forward
signature capture on the prefiller; the scheduler fetch and routing
decision are sub-percent. The signature cache occupies $0.24\%$ of HBM
and each request carries a $12$~KiB signature, negligible against the
multi-megabyte KV transfer. This is why median TTFT stays within the
across-baseline spread in the main results.

\begin{table}[t]
\centering
\small
\setlength{\tabcolsep}{4pt}
\caption{\sys runtime overhead (Qwen3-30B-A3B, task, 8P16D, $60$~req/s;
median TTFT $69$~ms).}
\label{tab:overhead}
\vspace{-1.0em}
\begin{tabular}{llrr}
\toprule
Component & Locus & P50 & \% TTFT \\
\midrule
\texttt{record()} hook      & prefill host & $0.02$\,ms          & $<0.1$ \\
\texttt{reduce()} scatter   & prefill GPU  & $0.48$\,ms          & $0.7$ \\
\texttt{stage\_sigs()} D2H  & prefill GPU  & $0.21$\,ms          & $0.3$ \\
\texttt{pop\_sig} fetch     & scheduler    & $7.0$\,\textmu s    & $<0.1$ \\
Route ($\tau$-JSQ)          & proxy        & $0.15$\,ms          & $0.2$ \\
\midrule
\textbf{Total}              &              & $\mathbf{0.86}$\,\textbf{ms} & $\mathbf{1.2}$ \\
\bottomrule
\end{tabular}
\end{table}

\subsection{Design Validation}
\label{sec:design-validation}

\begin{figure}[t!]
\centering
\includegraphics[width=0.95\columnwidth]{figure/06_evaluation/fig_expert_reduction.pdf}
\vspace{-1.0em}
\caption{Mean active experts per decode step on Qwen3-30B-A3B in the task domain on
an 8P16D cluster.}
\label{fig:expert_reduction}
\end{figure}

\paragraphnew{Active-Expert Reduction}
\label{sec:expert-reduction}
We validate that \sys's TPOT improvements come from a reduction in
the number of distinct experts activated per decode step.
\Cref{fig:expert_reduction} measures this count on
Qwen3-30B-A3B serving the task domain at 8P16D, aggregated over
$\sim$50K decode steps. \sys reduces the per-step active-expert
count by $22.0\%$ on average across decode batch sizes compared
to RR.

\begin{figure}[t!]
\centering
\includegraphics[width=0.95\columnwidth]{figure/06_evaluation/fig_transform_ablation.pdf}
\vspace{-1.0em}
\caption{TPOT P50/P99 \%\,$\Delta$ vs.\ RR at $r{=}60$\,req/s (8P16D, six
(model, dataset) cells) for two signature transforms: the IDF-reweighted
top-$k$ count (\textbf{count$\cdot$idf}) vs.\ the continuous softmax
gate (\textbf{gate-prob}). Rest of the recipe
fixed (greedy mask, balanced $K$-means $K{=}16$, $\tau{=}0.1$).}
\label{fig:transform_ablation}
\end{figure}

\paragraphnew{Expert Signature}
We validate \sys's signature choice: the IDF-reweighted top-$k$ count
(count$\cdot$idf).
\Cref{fig:transform_ablation} compares it against the
strongest continuous alternative, the softmax gate-probability,
to validate the discrete signature design
(\cref{sec:design_signature}). count$\cdot$idf reduces TPOT
P50 by an additional $3$\,pp on average and up to $14$\,pp across
the six cells, confirming that $\rho$ faithfully ranks signature
quality and that discrete prefill counts are the right primitive.

\begin{figure}[t!]
\centering
\includegraphics[width=0.95\columnwidth]{figure/06_evaluation/fig_cluster_ablation.pdf}
\vspace{-1em}
\caption{Mean \%$\Delta$ vs.\ RR over five request rates
($20$--$100$\,qps) at 8P16D with $\tau{=}0.1$.}
\label{fig:cluster_ablation}
\end{figure}

\paragraphnew{Offline Cluster Balance}
\label{sec:eval_cluster_ablation}
\Cref{fig:cluster_ablation} validates \sys's choice of
Hungarian-balanced $K$-means over vanilla $K$-means at
$\tau{=}0.1$. Vanilla $K$-means reduces median TPOT by up to
$9.8\%$ but \emph{regresses} tail TPOT by as much as $17.4\%$:
the per-decoder locality win is real, but the load imbalance pushes
the tail well past round-robin. Hungarian-balanced $K$-means recovers
both metrics, reducing P50 by up to $12.6\%$ and P99 by up to
$6.8\%$---uniform decoder utilization keeps the tail in check
without sacrificing the median win.

\begin{figure}[t!]
\centering
\includegraphics[width=0.95\columnwidth]{figure/06_evaluation/fig_tau_ablation.pdf}
\vspace{-0.5em}
\caption{Mean \%$\Delta$ vs RR (five rates, 20--100 qps; 8P16D) for six
cells (three models $\times$ two datasets) at four $\tau$ values.}
\label{fig:tau_ablation}
\end{figure}

\paragraphnew{Locality Band Width}
\label{sec:eval_tau_ablation}
\Cref{fig:tau_ablation} validates \sys's choice of
$\tau{=}0.1$ across $\tau \in \{0, 0.1, 0.2, 0.3\}$. Pure top-1
routing ($\tau{=}0$) regresses tail TPOT on four of six
workloads---by as much as $7.9\%$ on Gemma language and $6.7\%$
on Gemma task---because a transient burst of similar requests
lands on one decoder. A small band of $\tau{=}0.1$ removes the
regression on every workload while reducing median TPOT by
$5.2$--$12.7\%$ relative to RR. Beyond that the tail reduction
saturates and median TPOT erodes as more requests spill outside
their locality band, so \sys settles on $\tau{=}0.1$.

\begin{figure}[t!]
\centering
\includegraphics[width=0.95\columnwidth]{figure/06_evaluation/fig_prefix_cache_gptoss_task.pdf}
\vspace{-1.0em}
\caption{Composition with prefix caching on GPT-OSS-120B in the task
domain on an 8P16D cluster at $r{=}100$\,req/s, with
requests sampled cyclically from 2000 task prompts.}
\label{fig:prefix_ablation}
\vspace{0.6em}
\captionof{table}{Topology generalization on Qwen3-30B-A3B, language workload:
mean \%$\Delta$ vs RR across 20--100~qps.}
\label{tab:topology}
\begin{tabular}{lccc}
\toprule
 & 8P8D & 8P16D & 8P24D \\
\midrule
TPOT P50 & $-8.0\%$ & $-9.8\%$ & $-10.2\%$ \\
TPOT P99 & $-2.5\%$ & $-0.8\%$ & $-1.1\%$ \\
\bottomrule
\end{tabular}
\end{figure}

\paragraphnew{Prefix Cache Composition}
\label{sec:eval_prefix_cache}
\Cref{fig:prefix_ablation} validates that \sys composes with
the prefix cache. Turning the cache on collapses TTFT for both
routers without changing TPOT, and \sys's TPOT advantage over RR
(${\approx}13\%$ on P50, $5\%$ on P99) is preserved across both
cache states. The locality benefit and the prefix-cache benefit are
additive. A secondary back-pressure effect shows up on TTFT: \sys's
faster decoders drain requests sooner, keeping prefill queues
shorter and reducing TTFT when the cache is off.

\subsection{Generalization}
\label{sec:eval_generalization}
\paragraphnew{Decoder Pool Size}
To check that \sys generalizes across decoder counts, we sweep three
prefill--decode topologies on Qwen3-30B-A3B (language workload,
$TP{=}1$), holding the prefiller pool at 8 and growing the decoder pool:
8P8D (16 GPUs), 8P16D (24), and 8P24D (32). Over five rates
(20--100~qps), the mean median-TPOT reduction vs.\ round-robin grows
monotonically with the pool---$8.0\%$, $9.8\%$, $10.2\%$
(\Cref{tab:topology})---while tail TPOT stays within noise of
round-robin throughout. This follows the mechanism: more decoders split
each workload into finer clusters with narrower expert coverage, so the
per-decoder expert union shrinks and the latency reduction scales with
it. \sys thus generalizes across topologies, improving as the pool
expands.

\paragraphnew{Large MoE with Expert Parallelism}
We scale \sys to Qwen3-235B-A22B at 2P8D, $TP{=}4$, $EP{=}4$ per instance
(40 GPUs, 5 nodes), where each decoder shards attention and experts
across four GPUs. Clustering alone is insufficient here---a cluster's hot
experts can concentrate on one EP rank and bottleneck the decoder---so
each \sys decoder pairs clustering with the per-decoder expert placement
of~\cref{sec:implementation}, spreading each cluster's expected expert
work evenly across the four GPUs. On WildChat (\Cref{fig:qwen235},
24--56~req/s), \sys reduces median TPOT by $2.7$--$4.3\%$ and tail TPOT
by $0.6$--$2.0\%$ at every rate, confirming that locality routing
generalizes to large-MoE deployments with expert parallelism.

\begin{figure}[t!]
\centering
\includegraphics[width=0.95\columnwidth]{figure/06_evaluation/fig_qwen235_2p8d.pdf}
\vspace{-1.0em}
\caption{TPOT (median, p99) for Qwen3-235B-A22B on language workload
at 2P8D with $TP{=}4$ and $EP{=}4$ per instance (40 GPUs across 5 nodes).}
\label{fig:qwen235}
\end{figure}


% ===== FILE: sources/2607.00466-eldr-expert-locality-aware/arxiv/text/07_related.tex =====
\section{Related Work}
\label{sec:related}

\paragraphnew{PD disaggregated serving}
Prefill/decode (PD) disaggregation runs compute-bound prefill and
memory-bandwidth-bound decode on separate workers, removing interference
and meeting phase SLOs~\cite{distserve,splitwise,mooncake}; we adopt its
\emph{xPyD} topology. Cache-aware routers exploit prefix/KV
locality~\cite{sglang,preble} but treat decode workers as interchangeable
for expert computation. \sys complements them by adding
\emph{expert-activation} locality as the decode-routing signal, leaving
outputs unchanged.

\paragraphnew{MoE Expert locality}
Expert parallelism (EP) shards experts across devices, where skewed
routing causes load imbalance. One line rebalances \emph{within} a
deployment, replicating hot experts and rerouting tokens across EP
ranks~\cite{eplb,metro,moetuner,leastloadedep}---\emph{intra}-worker
balancing
that evens per-GPU expert load but not which worker serves a request.
Closest to us, systems exploiting expert-activation
locality---clustering requests by prefill activations~\cite{eapatterns}
or colocating experts with their tokens~\cite{semanticparallelism}---target
inter-node all-to-all \emph{communication}, not decode bandwidth. A third
line hits the same bandwidth bottleneck but \emph{approximately}, altering
decode expert selection: dropping low-importance experts~\cite{lynx},
piggybacking on already-loaded ones~\cite{opportunistic}, or sharing
experts across a batch~\cite{xshare}, trading accuracy for speed.

\sys differs on both axes. It performs \emph{inter}-worker balancing,
routing whole requests \emph{across} decode workers by expert locality to
shrink each worker's active experts. It is also \emph{lossless}: it
changes only which worker serves a request, never a token's expert
selection, so outputs match standard top-$k$ gating. It is thus exact
where batch-aware methods approximate, and orthogonal to intra-worker
balancing (each decode worker can still run EPLB internally), composing
with these methods rather than replacing them.


% ===== FILE: sources/2607.00466-eldr-expert-locality-aware/arxiv/text/08_conclusion.tex =====
\section{Conclusion}
\label{sec:conclusion}
Decode routing in PD-disaggregated MoE serving optimizes only
decode-worker load. \sys adds a second axis, \emph{expert locality}: a
prefill-derived signature drives load-tolerant routing that stays
coherent with the prefix cache, losslessly lowering median and tail TPOT
across models, workloads, and topologies.
